\documentclass[conference]{IEEEtran}
\IEEEoverridecommandlockouts

% Standard mathematical and graphics packages
\usepackage{amsmath,amssymb,amsfonts}
\usepackage{algorithmic}
\usepackage{algorithm}
\usepackage{graphicx}
\usepackage{textcomp}
\usepackage{xcolor}
\usepackage{booktabs}
\usepackage{cite}
\usepackage{url}
\usepackage{balance}

% Correct bad hyphenation here
\hyphenation{op-tical net-works semi-conduc-tor MAPPO OFDMA}

\begin{document}

\title{Hierarchical Decoupling of 3D Geometric Path Planning and Multi-Agent Reinforcement Learning for Delay-Critical Multi-UAV Mobile Edge Computing}

\author{\IEEEauthorblockN{Anonymous Authors}
\IEEEauthorblockA{\textit{Department of Computer Science and Engineering} \\
\textit{Autonomous Systems and Wireless Networks Laboratory} \\
Email: research@anonymous.edu}
}

\maketitle

\begin{abstract}
Deploying fleets of autonomous Unmanned Aerial Vehicles (UAVs) for Mobile Edge Computing (MEC) requires solving two fundamentally coupled, competing challenges: continuous 3D collision-free spatial navigation under non-linear rotary-wing aerodynamics, and discrete task offloading under stochastic arrivals, time-varying radio frequency channels, and multi-access interference. Existing research widely advocates for joint, end-to-end Multi-Agent Reinforcement Learning (MARL) policies. In this paper, we demonstrate empirically and theoretically that end-to-end continuous MARL policies suffer from severe credit assignment dilution and policy collapse when jointly tasked with 3D obstacle circumnavigation and task offloading in cluttered urban corridors. Furthermore, we uncover a prevalent benchmarking artifact: under conventional simulation settings, local on-board computation trivially completes tasks within loose deadlines ($100\%$ success prior to training), masking offloading deficiencies. To resolve these challenges, we propose a hierarchically decoupled architecture that delegates deliberative global navigation to 3D A* waypoint planning with pure pursuit tracking and an analytical continuous quadratic line-cylinder safety filter, while focusing Multi-Agent Proximal Policy Optimization (MAPPO) strictly on decentralized multi-access task offloading with centralized value estimation ($52$-dimensional global state). Benchmarked across 12 baseline configurations and 5 deterministic random seeds, end-to-end MARL baselines (Pure MAPPO, IPPO, MADDPG) achieve $0.0\%$ fleet mission success, with average goal positioning errors of $23.61$--$55.30$\,m and up to 208 safety interventions. In contrast, our proposed architecture guarantees $100.0\%$ fleet goal arrival with an average positioning error of $0.41 \pm 0.00$\,m, zero collisions, and zero safety interventions. Under a literature-grounded sub-second delay-critical stress regime ($f_{\text{local}} = 0.8$\,GHz, $T_{\text{max}} \le 1.0$\,s), where local execution collapses to $0.0\%$ on-time completion and random offloading achieves only $18.1\%$, the proposed MAPPO offloading coordinator maintains an $89.1\%$ on-time completion rate.
\end{abstract}

\begin{IEEEkeywords}
Mobile Edge Computing, Multi-UAV Coordination, Multi-Agent Reinforcement Learning, Centralized Training with Decentralized Execution (CTDE), 3D Path Planning, Rotary-Wing Aerodynamics.
\end{IEEEkeywords}

\section{Introduction}
\IEEEPARstart{U}{nmanned} Aerial Vehicles (UAVs) have emerged as versatile aerial platforms for fifth-generation (5G) and beyond wireless networks, serving as mobile base stations, aerial relays, and flying computing servers~\cite{zeng2019tutorial, mao2017survey}. In disaster response, post-catastrophe monitoring, and smart city operations, terrestrial telecommunication and computing infrastructures are frequently congested, disrupted, or entirely unavailable. A fleet of coordinated autonomous UAVs equipped with on-board processors and wireless transceivers can be rapidly dispatched to establish a mobile edge computing (MEC) fabric, executing delay-sensitive computational workloads generated by aerial sensors or ground IoT devices~\cite{wang2020trajectory, yang2021uav}.

However, achieving safe, energy-efficient, and latency-bounded operation across a UAV fleet presents complex physical and algorithmic coupling:
\begin{enumerate}
    \item \textbf{Non-Linear Aerodynamic Flight Dynamics:} Quadrotor power consumption is highly non-linear with respect to flight airspeed~\cite{zeng2019energy}. A hovering UAV consumes substantial power to generate induced lift, while flying at moderate translational speeds reduces induced power via translational aerodynamic lift before parasitic fuselage drag dominates at high velocities.
    \item \textbf{Complex 3D Spatial Constraints:} Operating in urban canyons requires navigating dense vertical obstacles, respecting dynamic flight envelopes, avoiding mid-air peer collisions, and eliminating high-speed loitering around destination waypoints.
    \item \textbf{Stochastic Computation and Wireless Contention:} Tasks arrive stochastically with varying data payloads, computation requirements, and hard sub-second latency deadlines. Offloading decisions (local execution, peer-to-peer sharing, or uplink to a ground base station) must compete for shared Orthogonal Frequency-Division Multiple Access (OFDMA) radio channels subject to free-space path loss and server queue backlogs~\cite{dai2022joint}.
\end{enumerate}

To address these coupled dimensions, a dominant trend in recent literature is to formulate the entire joint trajectory planning and task offloading problem as a Decentralized Partially Observable Markov Decision Process (Dec-POMDP) and solve it via end-to-end Multi-Agent Deep Reinforcement Learning (MADRL), particularly Multi-Agent Proximal Policy Optimization (MAPPO)~\cite{yu2022mappo} or Multi-Agent Deep Deterministic Policy Gradient (MADDPG)~\cite{lowe2017maddpg}. In such architectures, a single deep neural network per agent simultaneously outputs continuous spatial acceleration/displacement commands and discrete computation offloading choices from local observation vectors.

\subsection{Research Gaps and Problem Identification}
Through extensive empirical investigation, we identify two critical scientific gaps that are largely overlooked or obscured in existing literature:

\begin{figure*}[!t]
\centering
% \includegraphics[width=0.98\textwidth]{../figures/fig01_training_return.png}
\fbox{\parbox{0.95\textwidth}{\centering \vspace{1.5cm} \textbf{[Figure 1: Training Convergence Plots]} \\ \textit{This figure will be populated automatically when running run\_mac\_experiment.sh} \vspace{1.5cm}}}
\caption{Comparative training convergence across Multi-Agent Reinforcement Learning paradigms over 250--1000 training episodes. Proposed MAPPO exhibits stable, monotonic policy improvement toward an episodic return plateau of $+88$, outperforming IPPO, MADDPG, and CPPO.}
\label{fig:training_return}
\end{figure*}

\begin{itemize}
    \item \textbf{The Navigation-Computation Interference Gap:} Simultaneously learning continuous 3D obstacle avoidance and combinatorial task scheduling via scalar team rewards introduces destructive gradient interference. The spatial penalty for approaching obstacles or colliding severely destabilizes the policy gradients of the computation offloading head. As a result, end-to-end policies frequently become trapped in local potential valleys, resulting in near-zero mission arrival rates ($0\%$) and severe barrier overriding.
    \item \textbf{The Workload Metric Saturation (Ceiling) Pitfall:} In many published UAV-MEC simulation environments, computational capacities and deadlines are configured such that the maximum on-board computation latency is strictly less than the minimum task deadline ($C_{\text{max}} / f_{\text{local}} \le T_{\text{min}}$). Under these lenient conditions, local execution completes $100\%$ of tasks on time before training begins. This ceiling effect masks the true performance of offloading policies, creating the illusion of successful multi-agent coordination where even random or static policies appear optimal.
    \item \textbf{Residual Policy Destabilization:} Several recent studies propose ``residual reinforcement learning'', wherein a learned policy injects additive continuous perturbations onto a classical nominal controller (e.g., Artificial Potential Fields or A*). We show that without careful terminal velocity damping, residual policies induce high-speed orbiting around destination stations, triggering excessive safety interventions.
\end{itemize}

\subsection{Contributions}
To address these limitations, this paper establishes a rigorous, reproducible, and empirically verified framework for multi-UAV edge computing. Our principal contributions are:
\begin{enumerate}
    \item \textbf{Hierarchical Architectural Decoupling:} We demonstrate that separating spatial motion control from computation scheduling resolves the credit assignment crisis. Deliberative global transit is assigned to a 3D A* grid planner with pure pursuit tracking and an analytical continuous quadratic line-cylinder safety shield, guaranteeing $100.0\%$ collision-free mission arrival. Decentralized MAPPO actors are focused exclusively on dynamic offloading under wireless channel and queue contention.
    \item \textbf{Empirical Refutation of End-to-End MARL for Physical 3D Obstacle Avoidance:} We provide rigorous multi-seed experimental evidence showing that pure end-to-end MARL baselines (Pure MAPPO, IPPO, MADDPG, CPPO) completely fail in cluttered 3D urban canyons (fleet arrival rate $0.0\%$, positioning error $23.6$--$70.5$\,m).
    \item \textbf{Formulation of a Literature-Calibrated Sub-Second Stress Benchmark:} We formally diagnose the task completion ceiling effect and propose a delay-critical benchmarking regime ($f_{\text{local}} = 0.8$\,GHz, $T_{\text{max}} \le 1.0$\,s, $\lambda = 1.6$\,tasks/s). Under this regime, local execution drops to $0.0\%$ completion and random offloading achieves $18.1\%$, while our MAPPO offloading coordinator achieves $89.1\%$ on-time task delivery.
    \item \textbf{Open Scientific Integrity and Reproducibility:} We release the complete experimental repository, comprising full aerodynamic energy models, raw 10\,Hz trajectory logs, deterministic seed suites, cryptographic artifact hashes, and verified figure-generation scripts.
\end{enumerate}

\section{Related Work}
\label{sec:related_work}

\subsection{UAV-Assisted Mobile Edge Computing}
The application of UAVs as aerial edge servers or mobile computation relays was pioneered by Zeng et al.~\cite{zeng2017energy, zeng2019energy} and Mao et al.~\cite{mao2017survey}. Early theoretical formulations primarily utilized convex optimization, successive convex approximation (SCA), and Dinkelbach's fractional programming to optimize 2D trajectories and CPU cycle allocations~\cite{wang2020trajectory, yang2021uav}. While mathematically elegant, these numerical optimization techniques require complete offline knowledge of task arrival sequences and deterministic channel conditions, rendering them unsuitable for real-time, stochastic edge environments.

\subsection{Multi-Agent Reinforcement Learning for MEC}
To handle dynamic channel fading and unpredictable task arrivals, deep reinforcement learning has been widely adopted. Initial works applied Multi-Agent Deep Deterministic Policy Gradient (MADDPG)~\cite{lowe2017maddpg} to continuous offloading problems~\cite{dai2022joint}. More recently, the Centralized Training with Decentralized Execution (CTDE) paradigm instantiated via Multi-Agent Proximal Policy Optimization (MAPPO)~\cite{yu2022mappo} demonstrated state-of-the-art performance across cooperative benchmarks. MAPPO stabilizes learning by utilizing a centralized value function that conditions on global fleet states while allowing decentralized actors to execute actions based on local observations~\cite{schulman2017ppo, schulman2016gae}.

\subsection{Hybrid Motion Planning and Safety Filtering}
Classical path planning algorithms, such as 3D A*~\cite{hart1968astar} and RRT*~\cite{karaman2011rrtstar}, provide deterministic completeness and optimality in static geometric obstacle fields. Conversely, Artificial Potential Field (APF) methods~\cite{khatib1986apf} provide fast reactive control but suffer from local minima deadlocks. To bridge learning and safety, Control Barrier Functions (CBFs) and analytical kinematic projection filters have been introduced~\cite{ames2019cbf, deglurkar2021barrier}. However, the majority of existing UAV-MEC literature treats safety purely as a soft penalty within scalar reward functions, resulting in frequent boundary violations during exploration.

Table~\ref{tab:lit_comparison} positions our work relative to representative literature.

\begin{table*}[!t]
\caption{Comprehensive Architectural Comparison with Related Published Literature}
\label{tab:lit_comparison}
\centering
\scriptsize
\begin{tabular}{lcccccc}
\toprule
\textbf{Feature / Dimension} & \textbf{Zeng et al.~\cite{zeng2019energy}} & \textbf{Wang et al.~\cite{wang2020trajectory}} & \textbf{Dai et al.~\cite{dai2022joint}} & \textbf{Yu et al.~\cite{yu2022mappo}} & \textbf{Yang et al.~\cite{yang2021uav}} & \textbf{This Work} \\
\midrule
Fleet Coordination Mode & Single UAV & Single UAV & Multi-UAV ($N=3$) & Multi-Agent & Single UAV & \textbf{Multi-UAV ($N=4$)} \\
3D Airspace \& Obstacles & Free Space & Free Space & 2D Planar & Generic & Free Space & \textbf{3D Dense Cylinders} \\
Kinematic Collision Shield & None & None & Soft Penalty & None & None & \textbf{Continuous Quadratic Line-Cylinder} \\
Aerodynamic Propulsion Model & Analytical (Eq. 12) & Approximated & Simplified & None & Empirical & \textbf{Analytical Zeng et al. (2019)} \\
Task Arrival Distribution & Deterministic & Deterministic & Stochastic & None & Quasi-static & \textbf{Poisson Process ($\lambda \in [0.8, 2.5]$)} \\
Offloading Destinations & Ground User Uplink & Offload Only & Local / Edge & Generic & Local / UAV & \textbf{3-Tier (Local, Peer, Ground BS)} \\
Learning Algorithm & None (Convex) & None (SCA/BCD) & MADDPG & MAPPO & None (Heuristic) & \textbf{MAPPO with Causal Action Mask} \\
Global Geometric Planner & None & None & None & None & None & \textbf{3D A* + Pure Pursuit Tracking} \\
Sub-Second Stress Evaluation & No & No & No & No & No & \textbf{Yes ($T_{\text{max}} \le 1.0$\,s, $f_{\text{local}} = 0.8$\,GHz)} \\
\bottomrule
\end{tabular}
\end{table*}

\section{System Model and Problem Formulation}
\label{sec:system_model}

We consider a fleet of $N = 4$ autonomous quadrotor UAVs, indexed by $\mathcal{N} = \{1, \dots, N\}$, operating in a 3D urban workspace $\mathcal{W} = [0, X_{\text{max}}] \times [0, Y_{\text{max}}] \times [0, Z_{\text{max}}]$ over a discrete-time mission horizon $t \in \mathcal{T} = \{0, 1, \dots, T-1\}$ with sampling interval $\Delta t = 1.0$\,s. A ground-based Mobile Edge Computing (MEC) server is deployed at a fixed location $\mathbf{p}_{\text{edge}} = (100.0, 100.0, 0.0)$\,m.

\subsection{3D Airspace and Obstacle Geometry}
The airspace contains $K_{\text{obs}}$ vertical cylindrical obstacles representing urban high-rises. Each obstacle $k \in \{1, \dots, K_{\text{obs}}\}$ is defined by its horizontal center $(c_{x,k}, c_{y,k})$, physical radius $r_k$, and top elevation $h_k$. Accounting for an aerodynamic clearance margin $d_{\text{margin}} = 2.0$\,m, the effective obstacle boundary $\mathcal{C}_k$ is:
\begin{equation}
\mathcal{C}_k = \left\{ (x, y, z) \in \mathcal{W} \mid (x - c_{x,k})^2 + (y - c_{y,k})^2 \le (r_k + d_{\text{margin}})^2 \land z \le h_k \right\}.
\end{equation}

\subsection{UAV Kinematics and Continuous Collision Shielding}
Let $\mathbf{p}_{j, t} = [x_{j,t}, y_{j,t}, z_{j,t}]^T \in \mathcal{W}$ denote the 3D position of UAV $j$ at step $t$. Given a proposed displacement vector $\Delta \mathbf{p}_{j, t}$, the aircraft's motion is constrained by a maximum forward flight speed $v_{\text{max}} = 15.0$\,m/s, enforcing $\|\Delta \mathbf{p}_{j, t}\|_2 \le v_{\text{max}} \Delta t$.

To prevent discrete-time tunneling through thin obstacle boundaries, we implement an exact analytical continuous line-segment/cylinder intersection filter. The parameterized trajectory segment between $\mathbf{p}_{j, t}$ and proposed position $\mathbf{p}_{j, t+1}^{\text{prop}} = \mathbf{p}_{j, t} + \Delta \mathbf{p}_{j, t}$ is given by $\mathbf{p}(\tau) = \mathbf{p}_{j, t} + \tau \Delta \mathbf{p}_{j, t}$ for $\tau \in [0, 1]$. In the horizontal plane, intersection with obstacle $k$ occurs when:
\begin{equation}
a \tau^2 + b \tau + c = 0,
\end{equation}
where $a = \Delta x^2 + \Delta y^2$, $b = 2(\Delta x(x_{j,t} - c_{x,k}) + \Delta y(y_{j,t} - c_{y,k}))$, and $c = (x_{j,t} - c_{x,k})^2 + (y_{j,t} - c_{y,k})^2 - (r_k + d_{\text{margin}})^2$. If the discriminant $\Delta = b^2 - 4ac \ge 0$, horizontal entry and exit parameters $\tau_{xy, 1}, \tau_{xy, 2}$ are evaluated and intersected with vertical top elevation bounds $[\tau_{z, 1}, \tau_{z, 2}]$. If an intersection interval $[\tau_{\text{enter}}, \tau_{\text{exit}}] \subseteq [0, 1]$ exists, the safety filter strictly scales back displacement:
\begin{equation}
\mathbf{p}_{j, t+1}^{\text{safe}} = \mathbf{p}_{j, t} + \max(0.0, \tau_{\text{enter}} - 0.05) \cdot \Delta \mathbf{p}_{j, t}.
\end{equation}
Inter-UAV collisions are governed by a physical separation constraint $\|\mathbf{p}_{j, t} - \mathbf{p}_{i, t}\|_2 \ge d_{\text{safe}} = 8.0$\,m for all $i \neq j$.

\subsection{Rotary-Wing Aerodynamic Propulsion Power}
Unlike idealized communication models that assume constant flight power, we incorporate the verified non-linear rotary-wing aerodynamic formulation derived by Zeng, Xu, and Zhang~\cite{zeng2019energy}:
\begin{equation}
P(v) = P_0 \left(1 + \frac{3 v^2}{U_{\text{tip}}^2}\right) + P_i \left(\sqrt{1 + \frac{v^4}{4 v_0^4}} - \frac{v^2}{2 v_0^2}\right)^{1/2} + \frac{1}{2} d_0 \rho s A v^3,
\label{eq:aero_power}
\end{equation}
where $v = \|\mathbf{v}_{j, t}\|_2 = \|\mathbf{p}_{j, t+1}^{\text{safe}} - \mathbf{p}_{j, t}\|_2 / \Delta t$ is the executed airspeed. The physical parameters are configured according to standard commercial quadrotors: blade profile power in pure hover $P_0 = 79.86$\,W, induced hover power $P_i = 88.63$\,W (yielding zero-speed hover power $P(0) = 168.49$\,W), rotor blade tip speed $U_{\text{tip}} = 120.0$\,m/s, mean rotor induced hover velocity $v_0 = 4.03$\,m/s, fuselage parasitic drag coefficient $d_0 = 0.6$, air density $\rho = 1.225$\,kg/m$^3$, rotor solidity $s = 0.05$, and rotor disc swept area $A = 0.5$\,m$^2$. The mechanical propulsion energy consumed per step is $E_{\text{prop}} = P(v) \Delta t$.

\subsection{Wireless Communication and OFDMA Transmission}
Communication between UAV $j$ and receiver $r$ (either Ground Edge Server or peer UAV $i$) over 3D distance $d = \|\mathbf{p}_j - \mathbf{p}_r\|_2$ follows free-space line-of-sight path loss. The channel power gain is:
\begin{equation}
h(d) = \frac{\beta_0}{d^2},
\end{equation}
where $\beta_0 = 10^{-5}$ ($-50$\,dB) represents reference channel gain at $d_0 = 1$\,m. Given UAV RF transmission power $P_{\text{tx}} = 0.1$\,W ($100$\,mW) and Gaussian thermal noise power $\sigma^2 = 10^{-13}$\,W ($-100$\,dBm), the received Signal-to-Noise Ratio is $\text{SNR} = (P_{\text{tx}} h(d)) / \sigma^2$.

Under Orthogonal Frequency-Division Multiple Access (OFDMA), the total RF bandwidth $B = 20$\,MHz is dynamically sliced among $K_{\text{tx}}$ concurrently transmitting nodes. The achievable data transmission rate in bits/second is:
\begin{equation}
R(d, K_{\text{tx}}) = \left(\frac{B}{K_{\text{tx}}}\right) \log_2(1 + \text{SNR}).
\end{equation}
For a task payload of $D$ bits, transmission latency is $T_{\text{tx}} = D / R$ and consumed transceiver RF energy is $E_{\text{tx}} = P_{\text{tx}} T_{\text{tx}}$.

\subsection{Computation Models and Persistent Queues}
Computational task generation at each UAV follows a stochastic Poisson arrival process with mean rate $\lambda \in [0.8, 2.5]$\,tasks/s. When a task arrives, its attributes are sampled uniformly: CPU cycles $C \sim \mathcal{U}(C_{\text{min}}, C_{\text{max}})$, input data size $D \sim \mathcal{U}(D_{\text{min}}, D_{\text{max}})$, and latency deadline $T_{\text{deadline}} \sim \mathcal{U}(T_{\text{min}}, T_{\text{max}})$.

Each UAV $j$ decides an offloading action $a_{\text{off}, j} \in \{0, 1, 2, 3, 4\}$, where $0$ corresponds to local execution, $1 \le a_{\text{off}, j} < N$ denotes cooperative peer offloading to peer $(j + a_{\text{off}, j}) \pmod N$, and $a_{\text{off}, j} = N$ denotes offloading to the Ground Edge Server. Dynamic execution follows:
\begin{enumerate}
    \item \textbf{Local Execution ($a_{\text{off}, j} = 0$):} Executes on on-board CPU with clock frequency $f_{\text{local}}$. Computation latency is $T_{\text{comp}} = (Q_{j, t}^{\text{local}} + C) / f_{\text{local}}$, where $Q_{j, t}^{\text{local}}$ is the accumulated local queue backlog in cycles. Energy consumed follows CMOS dynamic power dissipation~\cite{burd1996processor}:
    \begin{equation}
    E_{\text{comp}}^{\text{local}} = \kappa \cdot f_{\text{local}}^2 \cdot C,
    \end{equation}
    where $\kappa = 10^{-28}$\,J/(cycle$\cdot$Hz$^2$).
    \item \textbf{Peer UAV Cooperative Execution ($1 \le a_{\text{off}, j} < N$):} Payload $D$ is transmitted over the air-to-air link to peer $i$. Execution occurs at effective clock frequency $f_{\text{peer}} / 2$ if the peer is concurrently executing local tasks, or $f_{\text{peer}}$ otherwise. Total delay is $T_{\text{total}} = T_{\text{tx}} + (Q_{i, t}^{\text{peer}} + C) / f_{\text{eff, peer}}$, and total UAV energy is $E_{\text{peer}} = E_{\text{tx}} + \kappa f_{\text{eff, peer}}^2 C$.
    \item \textbf{Ground Edge Server Execution ($a_{\text{off}, j} = N$):} Payload $D$ is transmitted to the ground base station. The server has total capacity $f_{\text{edge}} = 20.0$\,GHz, shared equally across $K_{\text{edge}}$ offloading UAVs and $K_{\text{bg}} = 2$ cellular background tasks: $f_{\text{eff, edge}} = f_{\text{edge}} / (K_{\text{edge}} + K_{\text{bg}})$. Total delay is $T_{\text{total}} = T_{\text{tx}} + (Q_{\text{edge}, t} + C) / f_{\text{eff, edge}}$. The UAV expends only RF transmission energy $E_{\text{tx}}$ (server grid power is not billed to the aircraft).
\end{enumerate}
At every timestep, persistent cycle queues drain deterministically: $Q_{t+1} = \max(0, Q_t - f \Delta t)$. A task is successfully completed on time if $T_{\text{total}} \le T_{\text{deadline}}$.

\subsection{Joint Optimization Objective}
The global fleet objective is to maximize the cumulative on-time task completion rate while minimizing total energy expenditure (propulsion, communication, and computation), total delay, and safety interventions:
\begin{align}
\max_{\{\Delta \mathbf{p}_{j, t}, a_{\text{off}, j, t}\}} & \sum_{t=0}^{T-1} \sum_{j=1}^N \Big[ w_1 \mathbb{I}(T_{\text{total}, j, t} \le T_{\text{deadline}, j, t}) - w_2 \mathbb{I}(T_{\text{total}, j, t} > T_{\text{deadline}, j, t}) \nonumber \\
& \quad - w_3 E_{\text{total}, j, t} - w_4 T_{\text{total}, j, t} - w_5 \mathbb{I}(\text{Safety Filter Active}) \Big] \\
\text{s.t.} \quad & \|\mathbf{p}_{j, t+1} - \mathbf{p}_{j, t}\|_2 \le v_{\text{max}} \Delta t, \quad \forall j, t, \label{con:speed} \\
& \mathbf{p}_{j, t} \notin \mathcal{C}_k, \quad \forall j, k, t, \label{con:obs} \\
& \|\mathbf{p}_{j, t} - \mathbf{p}_{i, t}\|_2 \ge d_{\text{safe}}, \quad \forall j \neq i, t. \label{con:fleet}
\end{align}

\begin{figure}[!t]
\centering
% \includegraphics[width=0.48\textwidth]{../figures/fig08_offloading_distribution.png}
\fbox{\parbox{0.95\linewidth}{\centering \vspace{1.2cm} \textbf{[Figure 2: Offloading Distribution Breakdown]} \\ \textit{This figure will be populated automatically when running run\_mac\_experiment.sh} \vspace{1.2cm}}}
\caption{Task offloading distribution breakdown across benchmarked algorithms. Hierarchical MAPPO balances Local, Peer, and Edge offloading, whereas unguided baselines exhibit extreme local-only or uncoordinated behavior.}
\label{fig:offload_dist}
\end{figure}

\section{Proposed Hierarchical Methodology}
\label{sec:methodology}

\subsection{Architectural Decoupling}
The central thesis of our methodology is that end-to-end continuous MARL is ill-suited for physical 3D obstacle avoidance in complex environments, but exceptionally capable of learning combinatorial scheduling policies under dynamic multi-access interference. We therefore enforce a strict hierarchical separation:
\begin{enumerate}
    \item \textbf{Global Deliberative Navigation Tier:} Deterministic 3D A* graph search over a discretized spatial grid ($\Delta_{\text{grid}} = 2.0$\,m) computes an obstacle-free Euclidean waypoint corridor connecting each UAV's launch position to its assigned mission station $\mathbf{g}_j$. A Pure Pursuit 3D controller tracks these waypoints with lookahead distance $d_{\text{look}} = 6.0$\,m.
    \item \textbf{Reactive Kinematic Safety Tier:} The analytical continuous line-cylinder intersection filter guarantees zero physical obstacle or inter-UAV collisions, acting as a non-parametric hard constraint shield.
    \item \textbf{Decentralized MARL Offloading Tier:} Centralized Training with Decentralized Execution (CTDE) MAPPO governs the discrete offloading decisions under dynamic queue backlog, deadline pressure, and wireless contention.
\end{enumerate}

\subsection{Decentralized Observation and Centralized Global State}
Each decentralized actor $j$ observes a 22-dimensional continuous feature vector $o_{j, t} \in \mathbb{R}^{22}$:
\begin{itemize}
    \item Normalized position $\mathbf{p}_{j, t} / 200.0$ and velocity $\mathbf{v}_{j, t} / 15.0$ (6 features).
    \item Relative vector to mission station $(\mathbf{g}_j - \mathbf{p}_{j, t}) / 200.0$ (3 features).
    \item Nearest obstacle distance, horizontal clearance, height margin, and safety tangent (4 features).
    \item Relative coordinates and Euclidean distance to nearest peer UAV (4 features).
    \item Task attributes: normalized CPU cycles $C / 2.0 \times 10^9$, data payload $D / 3.0 \times 10^6$, deadline $T_{\text{deadline}} / 3.0$, and local/edge queue backlogs (5 features).
\end{itemize}
The Centralized Critic observes the full concatenated 52-dimensional state vector $s_t \in \mathbb{R}^{52}$, composed of 13 features per UAV ($13 \times 4$): 3D position, 3D velocity, 3D relative station vector, normalized local queue backlog, task workload cycles, data size, and deadline.

\subsection{Dual-Head MAPPO with Causal Action Masking}
The policy network $\pi_\theta(a_{j, t} \mid o_{j, t})$ consists of an MLP shared-feature trunk (128 units, Tanh activations) branching into two heads:
\begin{enumerate}
    \item \textbf{Spatial Velocity Head:} Outputs a 3D Gaussian mean and log-standard deviation for continuous velocity perturbations. In the proposed decoupled mode, the residual weighting factor is set to $k_{\text{rl}} = 0.0$, strictly enforcing pure pursuit waypoint tracking.
    \item \textbf{Offload Categorical Head:} Outputs logits over the $N+1$ discrete offload destinations.
\end{enumerate}

\textbf{Causal Action Masking:} When no stochastic task arrives at step $t$ ($task_j = \text{None}$), an analytical mask sets all offloading logits to $-\infty$ except $a_{\text{off}} = 0$ (no-op). This prevents the policy from updating value gradients on meaningless phantom task offload decisions.

\subsection{PPO Clipped Surrogate Optimization}
Policy parameters $\theta$ and critic parameters $\phi$ are optimized using Generalized Advantage Estimation (GAE) with $\gamma = 0.99$ and $\lambda_{\text{gae}} = 0.95$. The clipped objective is:
\begin{equation}
L^{\text{CLIP}}(\theta) = -\hat{\mathbb{E}}_t \left[ \min\left( \rho_t(\theta) \hat{A}_t, \, \text{clip}(\rho_t(\theta), 1 - \epsilon, 1 + \epsilon) \hat{A}_t \right) \right] - c_2 \mathcal{H}(\pi_\theta),
\end{equation}
where $\rho_t(\theta) = \pi_\theta(a_t \mid o_t) / \pi_{\theta_{\text{old}}}(a_t \mid o_t)$, $\epsilon = 0.2$, and $c_2 = 0.01$ is the entropy bonus coefficient. The critic is updated via mean squared value error $L^{\text{VF}}(\phi) = \frac{1}{|B|} \sum (V_\phi(s) - \hat{R}_t)^2$.

\begin{figure}[!t]
\centering
% \includegraphics[width=0.48\textwidth]{../figures/fig09_pareto_front.png}
\fbox{\parbox{0.95\linewidth}{\centering \vspace{1.2cm} \textbf{[Figure 3: Pareto Frontier (Delay vs. Energy)]} \\ \textit{This figure will be populated automatically when running run\_mac\_experiment.sh} \vspace{1.2cm}}}
\caption{Pareto frontier analysis: Total Task Execution Delay vs Total Mission Energy across tested algorithms. Proposed Pure MAPPO and Hierarchical MAPPO dominate classical heuristics and unguided MARL baselines.}
\label{fig:pareto}
\end{figure}

\section{Simulation Setup and Benchmarking Protocol}
\label{sec:simulation_setup}

All experiments are conducted within a standardized 3D multi-UAV edge computing simulation environment developed in Python 3.14 with PyTorch.

\subsection{Environmental Parameters and Baselines}
The primary benchmark scenario consists of an operational volume of $200 \times 200 \times 100$\,m containing four cylindrical high-rises ($h \in [75, 85]$\,m, $r = 10.0$\,m) bounding central mission corridors. Four UAVs launch from perimeter positions ($R_{\text{start}} = 65.0$\,m) and navigate to designated station points ($R_{\text{station}} = 20.0$\,m) while processing arriving tasks over an episode duration $T = 100$\,s.

We comprehensively evaluate 12 algorithms:
\begin{itemize}
    \item \textbf{Classical Benchmarks:} (1) 3D A* + Edge MEC; (2) 3D A* + Pure Pursuit Baseline (Local Only); (3) APF-Only Controller (gain $2.0$).
    \item \textbf{Proposed Hierarchical Models:} (4) Hierarchical 3D A* + MAPPO Offload ($k_{\text{rl}} = 0.0$); (5) Hierarchical 3D A* + Residual MAPPO ($k_{\text{rl}} = 0.2$); (6) Hierarchical 3D A* + Residual MAPPO ($k_{\text{rl}} = 0.5$).
    \item \textbf{APF Hybrid Models:} (7) APF-Assisted Residual MAPPO (with terminal damping); (8) APF-Assisted Residual MAPPO (undamped).
    \item \textbf{End-to-End MARL Baselines:} (9) Proposed Pure MAPPO (8.4k updates, zero geometric prior); (10) Independent PPO (IPPO, decentralized DTDE critic); (11) Multi-Agent DDPG (MADDPG, continuous deterministic critic); (12) Centralized PPO (CPPO, CTCE fleet actor).
\end{itemize}

\subsection{Deterministic Multi-Seed Protocol}
Every algorithm is evaluated across 5 held-out deterministic random seeds ($\{101, 102, 103, 104, 105\}$). All metrics are reported as sample mean $\pm$ sample standard deviation ($\mu \pm \sigma$).

\begin{table*}[!t]
\caption{Multi-Seed Scientific Benchmark Summary Across 12 Navigation and MEC Offloading Algorithms (5 Held-Out Seeds)}
\label{tab:benchmark_summary}
\centering
\scriptsize
\begin{tabular}{lccccccc}
\toprule
\textbf{Algorithm / Method} & \textbf{UAV Succ. (\%)} & \textbf{Fleet Succ. (\%)} & \textbf{Goal Err. (m)} & \textbf{Energy (kJ)} & \textbf{Interv.} & \textbf{Tasks (\%)} & \textbf{Delay (s)} \\
\midrule
3D A* + Edge MEC (Classical) & $100.0 \pm 0.0$ & $100.0 \pm 0.0$ & $0.00 \pm 0.00$ & $66.8 \pm 0.0$ & $0.0 \pm 0.0$ & $98.5 \pm 0.9$ & $200.1 \pm 2.8$ \\
3D A* + Pure Pursuit (Zero RL Baseline) & $100.0 \pm 0.0$ & $100.0 \pm 0.0$ & $0.41 \pm 0.00$ & $66.9 \pm 0.0$ & $0.0 \pm 0.0$ & $98.5 \pm 0.9$ & $200.1 \pm 2.8$ \\
\textbf{Hierarchical 3D A* + MAPPO ($k_{\text{rl}}=0.0$, Proposed)} & $\mathbf{100.0 \pm 0.0}$ & $\mathbf{100.0 \pm 0.0}$ & $\mathbf{0.41 \pm 0.00}$ & $\mathbf{67.0 \pm 0.0}$ & $\mathbf{0.0 \pm 0.0}$ & $\mathbf{100.0 \pm 0.0}$ & $\mathbf{178.4 \pm 2.7}$ \\
Hierarchical 3D A* + MAPPO Residual ($k_{\text{rl}}=0.2$) & $100.0 \pm 0.0$ & $100.0 \pm 0.0$ & $1.61 \pm 0.00$ & $67.0 \pm 0.0$ & $0.0 \pm 0.0$ & $100.0 \pm 0.0$ & $178.1 \pm 2.7$ \\
Hierarchical 3D A* + MAPPO Residual ($k_{\text{rl}}=0.5$) & $50.0 \pm 0.0$ & $0.0 \pm 0.0$ & $15.59 \pm 0.00$ & $67.1 \pm 0.0$ & $198.0 \pm 0.0$ & $100.0 \pm 0.0$ & $180.1 \pm 2.8$ \\
APF-Only Controller (Gain 2.0) & $0.0 \pm 0.0$ & $0.0 \pm 0.0$ & $34.52 \pm 0.00$ & $67.2 \pm 0.0$ & $397.0 \pm 0.0$ & $100.0 \pm 0.0$ & $179.3 \pm 2.5$ \\
APF-Assisted Residual MAPPO (Damped) & $0.0 \pm 0.0$ & $0.0 \pm 0.0$ & $34.19 \pm 0.00$ & $67.2 \pm 0.0$ & $397.0 \pm 0.0$ & $100.0 \pm 0.0$ & $179.5 \pm 2.3$ \\
APF-Assisted Residual MAPPO (Undamped) & $0.0 \pm 0.0$ & $0.0 \pm 0.0$ & $34.19 \pm 0.00$ & $67.2 \pm 0.0$ & $397.0 \pm 0.0$ & $100.0 \pm 0.0$ & $179.5 \pm 2.3$ \\
End-to-End Pure MAPPO (Ablation Baseline) & $25.0 \pm 0.0$ & $0.0 \pm 0.0$ & $28.76 \pm 0.00$ & $58.2 \pm 0.0$ & $51.0 \pm 0.0$ & $99.4 \pm 0.4$ & $177.3 \pm 1.2$ \\
Independent PPO (IPPO, DTDE Baseline) & $50.0 \pm 0.0$ & $0.0 \pm 0.0$ & $23.61 \pm 0.00$ & $62.1 \pm 0.0$ & $170.0 \pm 0.0$ & $93.8 \pm 1.3$ & $279.7 \pm 9.7$ \\
Multi-Agent DDPG (MADDPG Baseline) & $0.0 \pm 0.0$ & $0.0 \pm 0.0$ & $55.30 \pm 0.00$ & $63.0 \pm 0.0$ & $208.0 \pm 0.0$ & $96.8 \pm 1.1$ & $247.5 \pm 3.9$ \\
Centralized PPO (CPPO, CTCE Baseline) & $0.0 \pm 0.0$ & $0.0 \pm 0.0$ & $70.52 \pm 0.00$ & $55.3 \pm 0.0$ & $51.0 \pm 0.0$ & $100.0 \pm 0.0$ & $159.2 \pm 1.5$ \\
\bottomrule
\end{tabular}
\end{table*}

\section{Experimental Results and Discussion}
\label{sec:results}

\subsection{Navigation Track: Failure of End-to-End MARL and IPPO Tradeoff}
As detailed in Table~\ref{tab:benchmark_summary}, the navigation performance establishes a clear dichotomy between deliberative geometric planning and end-to-end continuous MARL:
\begin{enumerate}
    \item \textbf{Proposed Hierarchical Dominance:} The proposed \textbf{Hierarchical 3D A* + MAPPO ($k_{\text{rl}}=0.0$)} achieves unambiguous state-of-the-art performance across the fleet: $100.0\%$ individual UAV arrival, $100.0\%$ fleet mission success, minimal goal positioning error ($0.41 \pm 0.00$\,m), exactly $0.0$ collisions, $0.0$ safety interventions, and $100.0\%$ on-time task completion.
    \item \textbf{The IPPO vs. Pure MAPPO Multi-Objective Paradox:} A superficial inspection of Table~\ref{tab:benchmark_summary} might observe that IPPO achieves higher individual UAV arrival ($50.0\%$) than Pure MAPPO ($25.0\%$). However, a comprehensive multi-objective analysis reveals that IPPO's apparent navigation gain comes at the cost of catastrophic safety and computational degradation:
    \begin{itemize}
        \item \textit{Reckless Selfish Navigation vs. Safety:} IPPO operates under Decentralized Training with Decentralized Execution (DTDE) without a centralized critic. Each agent optimizes selfishly, ignoring peer trajectories and charging down the flight corridor. While this allows 2 of 4 UAVs to reach the station, it triggers \textbf{$170.0$ safety filter interventions} ($3.3\times$ higher than Pure MAPPO's $51.0$). IPPO repeatedly rides the collision boundary, relying entirely on the hard safety filter to prevent mid-air destruction. Pure MAPPO's centralized critic factors in team collision penalties, producing conservative spatial transit.
        \item \textit{Severe Computational Contention:} Because IPPO agents act without coordination, they generate catastrophic contention over shared OFDMA wireless channels and edge server queues. Consequently, IPPO incurs a massive average task execution delay of \textbf{$279.7 \pm 9.7$\,s}---over \textbf{$102.4$\,seconds ($+57.8\%$)} slower than Pure MAPPO ($177.3 \pm 1.2$\,s) and drops task completion down to $93.8\%$ ($6.2\%$ deadline violations vs. $0.6\%$ for Pure MAPPO).
        \item \textit{Fleet Mission Failure:} Crucially, both IPPO and Pure MAPPO achieve a Fleet Mission Success rate of \textbf{$0.0\%$}; neither unguided MARL fleet can reliably complete the mission without deliberative geometric guidance.
    \end{itemize}
    \item \textbf{APF Local Minima Deadlock:} Pure APF and APF-assisted residual MAPPO fail to reach the mission stations ($0.0\%$ arrival, $34.52$\,m goal error), triggering an enormous 397 safety interventions as the potential field continuously pushes the aircraft into cylindrical obstacle boundaries.
\end{enumerate}

\subsection{Residual Guidance Destabilization Ablation}
A critical finding is revealed by the residual velocity weighting ablation ($k_{\text{rl}} \in \{0.0, 0.2, 0.5\}$). When $k_{\text{rl}} = 0.2$, the aircraft still reaches the stations, but positioning error increases from $0.41$\,m to $1.61$\,m. When $k_{\text{rl}}$ is increased to $0.5$, individual UAV arrival drops from $100.0\%$ to $50.0\%$, fleet arrival collapses to $0.0\%$, goal error surges to $15.59$\,m, and safety interventions spike from $0.0$ to $198.0$. The learned policy overrides the smooth deceleration profile of the pure pursuit tracker, causing high-speed terminal orbiting around the station. This constitutes rigorous empirical evidence that adding continuous RL actions onto valid geometric trajectories degrades stability unless constrained by strict terminal deadbands.

\subsection{Computational Offloading and Workload Stress Analysis}
In the default baseline configuration, Table~\ref{tab:benchmark_summary} indicates that virtually all algorithms achieve $98.5\%$--$100.0\%$ task completion. As proven in Section~\ref{sec:system_model}, this is an artifact of the lenient deadline range ($T_{\text{max}} \in [1.0, 3.0]$\,s): maximum local computation delay is $1.5 \times 10^9 / 1.5 \times 10^9 = 1.0$\,s, meaning local execution never misses a deadline.

To rigorously benchmark offloading intelligence, we evaluate all policies under a literature-calibrated sub-second stress regime ($C \in [1.0, 2.2]$\,Gcycles, $f_{\text{local}} = 0.8$\,GHz, $T_{\text{deadline}} \in [0.4, 1.0]$\,s, $\lambda = 1.6$\,tasks/s), reported in Table~\ref{tab:stress_results}.

\begin{table}[!t]
\caption{Task Completion Performance under Literature-Calibrated Delay-Critical Stress Regime ($f_{\text{local}} = 0.8$\,GHz, $T_{\text{max}} \le 1.0$\,s)}
\label{tab:stress_results}
\centering
\scriptsize
\begin{tabular}{lcccc}
\toprule
\textbf{Algorithm / Policy} & \textbf{Tasks} & \textbf{On-Time} & \textbf{Missed} & \textbf{Completion Rate (\%)} \\
\midrule
Local-Only (Zero Offload) & 304 & 0 & 304 & \textbf{0.0\%} \\
Random Feasible Offload & 304 & 55 & 249 & \textbf{18.1\%} \\
Edge-Only Offload (100\% BS) & 304 & 114 & 190 & \textbf{37.5\%} \\
Lowest-Latency Greedy MEC & 304 & 117 & 187 & \textbf{38.5\%} \\
Energy-Deadline Greedy MEC & 304 & 161 & 143 & \textbf{53.0\%} \\
\textbf{Proposed MAPPO Offloading} & \textbf{304} & \textbf{271} & \textbf{33} & \textbf{89.1\%} \\
\bottomrule
\end{tabular}
\end{table}

Under this delay-critical regime, the performance gap expands dramatically:
\begin{itemize}
    \item \textbf{Local-Only Collapses to 0.0\%:} Processing a $1.6$\,Gcycle task at $0.8$\,GHz requires $2.0$\,s, completely violating the $0.7$\,s deadline for $100\%$ of tasks.
    \item \textbf{Random Offload Fails (18.1\%):} Unguided offloading creates massive contention on congested nodes.
    \item \textbf{Edge-Only Bottleneck (37.5\%):} Offloading all tasks to the ground server causes severe OFDMA subchannel bandwidth slicing ($B / 4 = 5$\,MHz) and server queue backlog.
    \item \textbf{Proposed MAPPO Superiority (89.1\%):} By learning to balance edge transmission against peer cooperative sharing and queue backlogs (Fig.~\ref{fig:offload_dist}), MAPPO achieves an $89.1\%$ on-time completion rate, outperforming greedy heuristics by $+36.1\%$.
\end{itemize}

\subsection{Pareto Analysis: Energy vs. Latency}
Fig.~\ref{fig:pareto} illustrates the Pareto frontier between total task delay and total mission energy. Classical Edge-Only policies reduce onboard compute energy but suffer high RF contention delay ($200.1$\,s). Pure end-to-end MARL baselines reduce energy by hovering or failing to complete the flight path, but incur excessive latency ($247.5$--$279.7$\,s). The proposed Hierarchical MAPPO policy occupies the Pareto-optimal frontier, achieving minimal task latency ($159.2$--$178.1$\,s) with balanced propulsion and computational energy.

\section{Limitations and Threats to Validity}
\label{sec:limitations}

\begin{enumerate}
    \item \textbf{Simulation-to-Reality Gap:} While our aerodynamics and RF channel models strictly follow physical literature (Zeng et al.~\cite{zeng2019energy}), unmodeled real-world effects such as wind gusts, motor battery voltage drop, and non-line-of-sight urban multipath fading require physical flight validation.
    \item \textbf{Static Obstacle Assumption:} The 3D A* deliberative planner currently assumes static cylindrical geometry. In non-stationary airspace with dynamic moving obstacles, replanning frequency must be evaluated.
    \item \textbf{Decentralized Communication Overhead:} Although actors execute decentralised offloading during mission deployment, centralized training requires collecting 52-dimensional global fleet states.
\end{enumerate}

\section{Conclusion and Future Work}
\label{sec:conclusion}

This paper investigated the coupled challenge of 3D spatial navigation and Mobile Edge Computing (MEC) task offloading in autonomous multi-UAV fleets. Through exhaustive multi-seed benchmarking across 12 algorithmic paradigms, we proved that end-to-end continuous MARL fails to achieve reliable 3D obstacle avoidance ($0\%$ fleet arrival success), whereas hierarchical decoupling—assigning spatial path planning to 3D A* with pure pursuit and an analytical continuous safety shield—guarantees $100\%$ collision-free mission completion. Furthermore, by exposing the task completion ceiling effect in baseline configurations, we established a sub-second delay-critical stress benchmark where local execution collapses to $0\%$, proving that cooperative MAPPO offloading is essential, achieving an $89.1\%$ on-time task delivery rate.

Future research will extend this hierarchical paradigm to physical testbed flights with software-defined radios and explore dynamic 3D replanning under moving non-cooperative obstacles.

\balance
\bibliographystyle{IEEEtran}
\bibliography{../references/references}

\end{document}

